\uuid{i7Xp}
\titre{Convolution 2D à la main — mode valid}
\niveau{L2}
\module{Informatique / Signal}
\chapitre{Réseaux de neurones}
\sousChapitre{Convolution 2D}
\theme{Calcul à la main, noyau $2\times2$, mode valid}
\auteur{}
\datecreate{2026-03-23}
\organisation{}
\difficulte{1}

\contenu{
\texte{
  On considère l'image $I$ et le noyau $K$ suivants :
  \[
  I = \begin{pmatrix} 1 & 2 & 3 \\ 4 & 5 & 6 \\ 7 & 8 & 9 \end{pmatrix},
  \qquad
  K = \begin{pmatrix} 0 & 1 \\ 1 & 0 \end{pmatrix}.
  \]
  On utilise le mode \textbf{valid} (pas de padding) et un stride de $1$.
}
\begin{enumerate}
\item
  \question{Quelle est la taille de la carte de sortie ? Justifier à l'aide de la formule
  $\text{out} = \lfloor (\text{in} + 2p - k)/s \rfloor + 1$.}
  \reponse{
    Avec $\text{in} = 3$, $p = 0$, $k = 2$, $s = 1$ :
    \[
      \text{out} = \lfloor (3 + 0 - 2)/1 \rfloor + 1 = 1 + 1 = 2.
    \]
    La sortie est donc de taille $2 \times 2$.
  }

\item
  \question{Calculer à la main tous les coefficients de la carte de sortie $S = I \star K$.}
  \indication{Rappel : on renverse le noyau ($\mathrm{flip}$) avant de calculer le produit scalaire avec chaque fenêtre.}
  \reponse{
    Le noyau renversé est $K' = \mathrm{flip}(K) = \begin{pmatrix} 0 & 1 \\ 1 & 0 \end{pmatrix}$
    (symétrique, donc $K' = K$).

    Les quatre fenêtres $2 \times 2$ extraites de $I$ (en mode valid) sont :
    \begin{align*}
      F_{00} &= \begin{pmatrix}1&2\\4&5\end{pmatrix}, &
      F_{01} &= \begin{pmatrix}2&3\\5&6\end{pmatrix}, \\
      F_{10} &= \begin{pmatrix}4&5\\7&8\end{pmatrix}, &
      F_{11} &= \begin{pmatrix}5&6\\8&9\end{pmatrix}.
    \end{align*}
    Produits scalaires terme à terme avec $K'$ :
    \begin{align*}
      S_{00} &= 1\cdot0 + 2\cdot1 + 4\cdot1 + 5\cdot0 = 6,\\
      S_{01} &= 2\cdot0 + 3\cdot1 + 5\cdot1 + 6\cdot0 = 8,\\
      S_{10} &= 4\cdot0 + 5\cdot1 + 7\cdot1 + 8\cdot0 = 12,\\
      S_{11} &= 5\cdot0 + 6\cdot1 + 8\cdot1 + 9\cdot0 = 14.
    \end{align*}
    Donc :
    \[
      S = \begin{pmatrix} 6 & 8 \\ 12 & 14 \end{pmatrix}.
    \]
  }

\item
  \question{Que détecte intuitivement ce noyau $K$ ?}
  \reponse{
    Le noyau $K = \begin{pmatrix}0&1\\1&0\end{pmatrix}$ somme les deux pixels situés sur la
    diagonale anti-principale (haut-droite et bas-gauche) de chaque fenêtre.
    Il réagit aux variations diagonales (à $45°$) dans l'image, et peut être vu comme
    un détecteur de corrélation diagonale. Ce n'est ni un détecteur horizontal ni vertical pur.
  }
\end{enumerate}
}
